\subsubsection{The one-gluon incoming state}\label{sec:singleglue}
We now compute the dressed one-gluon state, $\ket{\Psi^{a}_{i}(k)}=\Omega\ket{g^{a}_{i}(k)}$, to order
$g^{2}$. Its components are organised by the number of outgoing soft gluons and the power of $\rho$. The
LCPT matrix elements are given in Eqs.~\eqref{m1}--\eqref{m5}, and the transverse Fourier integrals in
Eqs.~\eqref{i1}--\eqref{i5}. In transverse coordinates, $\v z'$ is the position of the incoming gluon,
$\v z$ that of an outgoing gluon, and $\v x,\v y$ are positions of valence charges. We write
\begin{align}\label{eq:XXY}
X\equiv\v x-\v z\,,\qquad X'\equiv\v x-\v z'\,,\qquad Y\equiv\v y-\v z\,,
\end{align}
and, for the outgoing gluon, $\xi=p^{+}/k^{+}$. The tensor
\begin{align}\label{eq:Ttensor}
\mathcal{T}_{ijlm}(\xi)\equiv\delta_{il}\,\delta_{jm}-\frac{1}{\xi}\,\delta_{jl}\,\delta_{im}-\frac{1}{1-\xi}\,\delta_{ij}\,\delta_{lm}
\end{align}
collects the polarisation structure of the triple-gluon vertex.

%%%%%%%%%%%%%%%%%%%%%%%%%%%%%%%%%%%%%%%%%%%%%%%%%%%%%%%%%%%%%%%%%%%%%%%%%%%%%%%%
\paragraph{No outgoing gluon.}
At order $g$ only one diagram contributes, Fig.~\ref{fig:gtovac}: the incoming gluon is absorbed by the
valence charges. Using the matrix element~\eqref{m1},
\begin{align}\label{eq:psirho}
\ket{\big(\Psi^{a}_{i}(k)\big)_{\rho}}=\ket{0}\,\frac{\bra{0}H_{g}\ket{g^{a}_{i}(k)}}{E_{g}(k)}
=\ket{0}\,\frac{g\,\v k^{i}}{2\pi^{3/2}\sqrt{k^{+}}\,\v k^{2}}\,\rho^{a}(\v k)\,.
\end{align}
In transverse coordinates, with Eq.~\eqref{i1},
\begin{align}\label{eq:psirhoX}
\ket{\big(\Psi^{a}_{i}(k^{+},\v z')\big)_{\rho}}=\ket{0}\,\frac{ig}{2\pi^{3/2}\sqrt{k^{+}}}\int_{\v x}\frac{X'^{i}}{X'^{2}}\,\rho^{a}(\v x)\,.
\end{align}
Since $(2\pi)^{-3/2}\mf{A}^{(1)}_{1}(k)$ is the amplitude of $\ket{0}$ in Eq.~\eqref{eq:Omegag}, this gives
$\mf{A}^{(1)a}_{1i}(k)=\sqrt{2}\,g\,\v k^{i}\rho^{a}(\v k)/(\sqrt{k^{+}}\,\v k^{2})=-\mf{A}^{(1)\dagger a}_{2i}(k)$, in
agreement with Eqs.~\eqref{eq:A1rel} and~\eqref{eq:A1N2explicit}.

\begin{figure}[t]
\centering
\tikzfig{g_tovacum}
\caption{Absorption of the incoming gluon by the valence charges. The vertical dashed line marks the
intermediate state.}
\label{fig:gtovac}
\end{figure}

%%%%%%%%%%%%%%%%%%%%%%%%%%%%%%%%%%%%%%%%%%%%%%%%%%%%%%%%%%%%%%%%%%%%%%%%%%%%%%%%
\paragraph{One outgoing gluon, linear in $\rho$.}
At order $g^{2}$ there are three diagrams with one outgoing gluon and one power of $\rho$,
Fig.~\ref{fig:onegrho}: the instantaneous interaction, (a), and the two time orderings of a valence
emission and the triple-gluon vertex, (b) and (c). In each case the plus momentum flows either from the
valence charges to the gluon or from the gluon to the valence charges. The two directions give
different ranges of $\xi$, fixed by the kinematic constraints of Fig.~\ref{fig:kinem}: $\xi>1$ when the
valence charges give plus momentum to the gluon, and $\xi<1$ when they take it.

\begin{figure}[t]
\centering
\begin{minipage}[b]{0.32\textwidth}\centering\tikzfig{gtog_inst}\\[2pt](a)\end{minipage}\hfill
\begin{minipage}[b]{0.32\textwidth}\centering\tikzfig{gandgtog}\\[2pt](b)\end{minipage}\hfill
\begin{minipage}[b]{0.32\textwidth}\centering\tikzfig{gto3g_torho}\\[2pt](c)\end{minipage}
\caption{One outgoing gluon, linear in $\rho$: (a) instantaneous interaction; (b) the valence charges
emit a gluon, which merges with the incoming one through the triple-gluon vertex; (c) the incoming
gluon splits through the triple-gluon vertex and one of the two gluons is absorbed by the valence
charges.}
\label{fig:onegrho}
\end{figure}

\begin{figure}[t]
\centering
\tikzfig{kinem}
\caption{Kinematic constraint on the plus momentum of the soft gluon. The allowed range of $p^{+}$ is
the intersection of the shaded regions. The red region is due to the constraint
$\Lambda<p^{+}-k^{+}<\vee$, the green region to $\Lambda<p^{+}<\vee$ and $\Lambda<k^{+}<\vee$.}
\label{fig:kinem}
\end{figure}

\emph{Instantaneous interaction}, Fig.~\ref{fig:onegrho}(a). When the plus momentum flows from the valence
charges to the gluon, $k^{+}+\Lambda<p^{+}<\vee$, and with the matrix element~\eqref{m5}
\begin{align}\label{oneginst1}
\ket{\big(\Psi^{a}_{i}(k)\big)^{1}_{g\rho,\,{\rm inst}}}
=\ &\int_{k^{+}+\Lambda}^{\vee}dp^{+}\int d^{2}\v p\,\ket{g^{b}_{j}(p)}\,
\frac{\bra{g^{b}_{j}(p)}H_{gg\text{-inst}}\ket{g^{a}_{i}(k)}}{E_{g}(k)-E_{g}(p)}\notag\\
=\ &\frac{ig^{2}f^{abc}}{(2\pi)^{3}}\int_{k^{+}+\Lambda}^{\vee}dp^{+}\int d^{2}\v p\,\ket{g^{b}_{j}(p^{+},\v p)}\,
\delta_{ij}\,\frac{\sqrt{p^{+}k^{+}}\,(k^{+}+p^{+})}{(p^{+}-k^{+})^{2}}\,
\frac{\rho^{c}(\v k-\v p)}{p^{+}\v k^{2}-k^{+}\v p^{2}}\,.
\end{align}
In terms of $\xi$,
\begin{align}\label{oneginst11}
\ket{\big(\Psi^{a}_{i}(k)\big)^{1}_{g\rho,\,{\rm inst}}}
=\frac{ig^{2}f^{abc}}{(2\pi)^{3}}\int_{1+\Lambda/k^{+}}^{\vee/k^{+}}d\xi\int d^{2}\v p\,
\ket{g^{b}_{j}(\xi k^{+},\v p)}\,\delta_{ij}\,\frac{\sqrt{\xi}\,(1+\xi)}{(1-\xi)^{2}}\,
\frac{\rho^{c}(\v k-\v p)}{\xi\v k^{2}-\v p^{2}}\,,
\end{align}
and in transverse coordinates, with Eq.~\eqref{i3},
\begin{align}\label{Eq:onegluoninstX1}
\ket{\big(\Psi^{a}_{i}(k^{+},\v z')\big)^{1}_{g\rho,\,{\rm inst}}}
=\frac{ig^{2}f^{abc}}{(2\pi)^{3}}\int_{1+\Lambda/k^{+}}^{\vee/k^{+}}d\xi\int_{\v x,\v z}
\ket{g^{b}_{j}(\xi k^{+},\v z)}\,\delta_{ij}\,\frac{\sqrt{\xi}\,(1+\xi)}{(1-\xi)^{2}}\,
\frac{\rho^{c}(\v x)}{\xi X^{2}-X'^{2}}\,.
\end{align}
The range of $\xi$ is unusual here, since $\xi>1$. When the plus momentum flows from the gluon to the
valence charges we obtain the same expression with $\Lambda/k^{+}<\xi<1-\Lambda/k^{+}$:
\begin{align}\label{oneginst2}
\ket{\big(\Psi^{a}_{i}(k)\big)^{2}_{g\rho,\,{\rm inst}}}
=\frac{ig^{2}f^{abc}}{(2\pi)^{3}}\int_{\Lambda/k^{+}}^{1-\Lambda/k^{+}}d\xi\int d^{2}\v p\,
\ket{g^{b}_{j}(\xi k^{+},\v p)}\,\delta_{ij}\,\frac{\sqrt{\xi}\,(1+\xi)}{(1-\xi)^{2}}\,
\frac{\rho^{c}(\v k-\v p)}{\xi\v k^{2}-\v p^{2}}\,,
\end{align}
\begin{align}\label{Eq:onegluoninstX2}
\ket{\big(\Psi^{a}_{i}(k^{+},\v z')\big)^{2}_{g\rho,\,{\rm inst}}}
=\frac{ig^{2}f^{abc}}{(2\pi)^{3}}\int_{\Lambda/k^{+}}^{1-\Lambda/k^{+}}d\xi\int_{\v x,\v z}
\ket{g^{b}_{j}(\xi k^{+},\v z)}\,\delta_{ij}\,\frac{\sqrt{\xi}\,(1+\xi)}{(1-\xi)^{2}}\,
\frac{\rho^{c}(\v x)}{\xi X^{2}-X'^{2}}\,.
\end{align}

\emph{Valence emission followed by the triple-gluon vertex}, Fig.~\ref{fig:onegrho}(b). Here the valence
charges give plus momentum to the gluon. With the matrix elements~\eqref{m2} and~\eqref{m3},
\begin{align}\label{onegrho1}
\ket{\big(\Psi^{a}_{i}(k)\big)^{1}_{g\rho}}
=\ &\int d^{3}p\,d^{3}q\,d^{3}r\,\ket{g^{b}_{j}(p)}\,
\frac{\bra{g^{b}_{j}(p)}H_{ggg}\ket{g^{c}_{l}(q)\,g^{d}_{m}(r)}\bra{g^{c}_{l}(q)\,g^{d}_{m}(r)}H_{g}\ket{g^{a}_{i}(k)}}
{2\big(E_{g}(k)-E_{gg}(q,r)\big)\big(E_{g}(k)-E_{g}(p)\big)}\\
% CHECK: the overall sign below was flipped (it read -ig^2 f/8pi^3) so that this expression reduces to the
% xi form of Eq. (onegrho12xi). Only the xi form is consistent with the cancellation of the 1/(1-xi)^2
% pinch against Eq. (oneginst11). Please check the sign of the matrix elements (m2), (m3) entering here.
=\ &\frac{ig^{2}f^{abc}}{8\pi^{3}}\int_{k^{+}+\Lambda}^{\vee}dp^{+}\int d^{2}\v p\,\ket{g^{b}_{j}(p^{+},\v p)}\,
\frac{(\v p-\v k)^{l}\,\rho^{c}(\v k-\v p)}{\sqrt{p^{+}(p^{+}-k^{+})k^{+}}\,(p^{+}-k^{+})^{3/2}}\,
\frac{k^{+}-p^{+}}{\big(\frac{\v k^{2}}{k^{+}}-\frac{\v p^{2}}{p^{+}}\big)(\v k-\v p)^{2}}\notag\\
&\times\bigg[\Big(2\v k^{j}-\v p^{j}+\frac{p^{+}-2k^{+}}{p^{+}}\,\v p^{j}\Big)\delta_{il}
+\Big(2\v p^{i}-\v k^{i}-\frac{2p^{+}-k^{+}}{k^{+}}\,\v k^{i}\Big)\delta_{jl}\notag\\
&\hspace{1cm}+\Big(\frac{p^{+}+k^{+}}{p^{+}-k^{+}}\,(\v p-\v k)^{l}-\v k^{l}-\v p^{l}\Big)\delta_{ij}\bigg]\,.\label{onegrho12}
\end{align}
In terms of $\xi$ the square bracket, contracted with $(\v p-\v k)^{l}$, reduces to
$2(\v p-\v k)^{m}(\v p-\xi\v k)^{l}\,\mathcal{T}_{ijlm}(\xi)$, and
\begin{align}\label{onegrho12xi}
\ket{\big(\Psi^{a}_{i}(k)\big)^{1}_{g\rho}}
=\ &-\frac{ig^{2}f^{abc}}{4\pi^{3}}\int_{1+\Lambda/k^{+}}^{\vee/k^{+}}d\xi\int d^{2}\v p\,
\ket{g^{b}_{j}(\xi k^{+},\v p)}\,\frac{\sqrt{\xi}}{1-\xi}\,
\frac{(\v p-\v k)^{m}\,(\v p-\xi\v k)^{l}}{(\v p-\v k)^{2}\,(\v p^{2}-\xi\v k^{2})}\,
\rho^{c}(\v k-\v p)\,\mathcal{T}_{ijlm}(\xi)\,.
\end{align}
With Eq.~\eqref{i4} the Fourier transform is
% CHECK: sign changed from +i to -i. The transform of (p-k)^m (p-xi k)^l/((p-k)^2 (p^2-xi k^2)) is
% +(X'-xi X)^m (X'-X)^l/((X'^2-xi X^2)(X'-X)^2) (checks/ft_onegluon.py), and only this sign makes the
% 1/(1-xi)^2 pinch cancel against Eq. (Eq:onegluoninstX1) in coordinate space.
\begin{align}\label{Eq:onegrhoX1}
\ket{\big(\Psi^{a}_{i}(k^{+},\v z')\big)^{1}_{g\rho}}
=\ &-\frac{ig^{2}f^{abc}}{4\pi^{3}}\int_{1+\Lambda/k^{+}}^{\vee/k^{+}}d\xi\int_{\v x,\v z}
\ket{g^{b}_{j}(\xi k^{+},\v z)}\,\frac{\sqrt{\xi}}{1-\xi}\,
\frac{X'^{m}-\xi X^{m}}{X'^{2}-\xi X^{2}}\,\frac{X'^{l}-X^{l}}{(X'-X)^{2}}\,
\rho^{c}(\v x)\,\mathcal{T}_{ijlm}(\xi)\,.
\end{align}

\emph{Triple-gluon vertex followed by valence absorption}, Fig.~\ref{fig:onegrho}(c). Here the gluon
gives plus momentum to the valence charges, and
\begin{align}\label{onegrho2}
\ket{\big(\Psi^{a}_{i}(k)\big)^{2}_{g\rho}}
=\ &\int d^{3}p\,d^{3}q\,d^{3}r\,\ket{g^{b}_{j}(p)}\,
\frac{\bra{g^{b}_{j}(p)}H_{g}\ket{g^{c}_{l}(q)\,g^{d}_{m}(r)}\bra{g^{c}_{l}(q)\,g^{d}_{m}(r)}H_{ggg}\ket{g^{a}_{i}(k)}}
{2\big(E_{g}(k)-E_{gg}(q,r)\big)\big(E_{g}(k)-E_{g}(p)\big)}\\
=\ &\frac{ig^{2}f^{abc}}{8\pi^{3}}\int_{\Lambda}^{k^{+}-\Lambda}dp^{+}\int d^{2}\v p\,\ket{g^{b}_{j}(p^{+},\v p)}\,
\frac{(\v k-\v p)^{l}\,\rho^{c}(\v k-\v p)}{\sqrt{k^{+}(k^{+}-p^{+})p^{+}}\,(k^{+}-p^{+})^{3/2}}\notag\\
&\times\frac{1}{\Big(\frac{\v k^{2}}{k^{+}}-\frac{(\v k-\v p)^{2}}{k^{+}-p^{+}}-\frac{\v p^{2}}{p^{+}}\Big)
\Big(\frac{\v k^{2}}{k^{+}}-\frac{\v p^{2}}{p^{+}}\Big)}\notag\\
&\times\bigg[\Big(2\v p^{i}-\v k^{i}+\frac{k^{+}-2p^{+}}{k^{+}}\,\v k^{i}\Big)\delta_{jl}
+\Big(2\v k^{j}-\v p^{j}-\frac{2k^{+}-p^{+}}{p^{+}}\,\v p^{j}\Big)\delta_{il}\notag\\
&\hspace{1cm}+\Big(\frac{k^{+}+p^{+}}{k^{+}-p^{+}}\,(\v k-\v p)^{l}-\v p^{l}-\v k^{l}\Big)\delta_{ij}\bigg]\,.\label{onegrho22}
\end{align}
In terms of $\xi$,
\begin{align}\label{onegrho22xi}
\ket{\big(\Psi^{a}_{i}(k)\big)^{2}_{g\rho}}
=\ &\frac{ig^{2}f^{abc}}{4\pi^{3}}\int_{\Lambda/k^{+}}^{1-\Lambda/k^{+}}d\xi\int d^{2}\v p\,
\ket{g^{b}_{j}(\xi k^{+},\v p)}\,\frac{\xi^{3/2}}{1-\xi}\,
\frac{(\v k-\v p)^{m}\,(\v p-\xi\v k)^{l}}{(\v p^{2}-\xi\v k^{2})\,(\v p-\xi\v k)^{2}}\,
\rho^{c}(\v k-\v p)\,\mathcal{T}_{ijlm}(\xi)\,,
\end{align}
and, with Eq.~\eqref{i5},
\begin{align}\label{Eq:onegrhoX2}
\ket{\big(\Psi^{a}_{i}(k^{+},\v z')\big)^{2}_{g\rho}}
=\ &-\frac{ig^{2}f^{abc}}{4\pi^{3}}\int_{\Lambda/k^{+}}^{1-\Lambda/k^{+}}d\xi\int_{\v x,\v z}
\ket{g^{b}_{j}(\xi k^{+},\v z)}\,\frac{\xi^{3/2}}{1-\xi}\,
\frac{X'^{l}-X^{l}}{X'^{2}-\xi X^{2}}\,\frac{X'^{m}-\xi X^{m}}{(X'-\xi X)^{2}}\,
\rho^{c}(\v x)\,\mathcal{T}_{ijlm}(\xi)\,.
\end{align}

The instantaneous terms contain $1/(1-\xi)^{2}$. The point $\xi=1$ lies outside the integration
ranges, but the limits $1\pm\Lambda/k^{+}$ pinch it, which could produce contributions proportional to
$k^{+}/\Lambda$. Near $\xi=1$ the $\delta_{ij}\delta_{lm}$ part of $\mathcal{T}_{ijlm}$ gives the
non-instantaneous terms the same double pole with the opposite sign. The singularity therefore cancels
between Eqs.~\eqref{Eq:onegluoninstX1} and~\eqref{Eq:onegrhoX1} and between
Eqs.~\eqref{Eq:onegluoninstX2} and~\eqref{Eq:onegrhoX2}, and likewise in momentum space.

%%%%%%%%%%%%%%%%%%%%%%%%%%%%%%%%%%%%%%%%%%%%%%%%%%%%%%%%%%%%%%%%%%%%%%%%%%%%%%%%
\paragraph{One outgoing gluon, quadratic in $\rho$.}
Two diagrams contribute, Fig.~\ref{fig:onegrhorho}. We remove the contributions in which the outgoing
gluon is identical to the incoming one. In the first diagram the incoming gluon is absorbed and a new
gluon is then emitted. With the matrix element~\eqref{m1},
\begin{align}\label{m1grr}
\ket{\big(\Psi^{a}_{i}(k)\big)^{1}_{g\rho\rho}}
=\ &\int_{\Lambda}^{\vee}dp^{+}\int d^{2}\v p\,\ket{g^{b}_{j}(p^{+},\v p)}\,
\frac{\bra{g^{b}_{j}(p)}H_{g}\ket{0}\bra{0}H_{g}\ket{g^{a}_{i}(k)}}{E_{g}(k)\big(E_{g}(k)-E_{g}(p)\big)}\notag\\
=\ &\frac{g^{2}}{4\pi^{3}}\int_{\Lambda/k^{+}}^{\vee/k^{+}}d\xi\int d^{2}\v p\,\ket{g^{b}_{j}(\xi k^{+},\v p)}\,
\frac{\v k^{i}\,\v p^{j}}{\sqrt{\xi}\,(\xi\v k^{2}-\v p^{2})\,\v k^{2}}\,\rho^{b}(-\v p)\,\rho^{a}(\v k)\,,
\end{align}
and, with Eq.~\eqref{i2},
\begin{align}\label{m1grrX}
\ket{\big(\Psi^{a}_{i}(k^{+},\v z')\big)^{1}_{g\rho\rho}}
=\frac{g^{2}}{4\pi^{3}}\int_{\Lambda/k^{+}}^{\vee/k^{+}}d\xi\int_{\v x,\v y,\v z}\ket{g^{b}_{j}(\xi k^{+},\v z)}\,
\frac{X'^{i}\,Y^{j}}{\sqrt{\xi}\,(\xi Y^{2}-X'^{2})\,Y^{2}}\,\rho^{b}(\v y)\,\rho^{a}(\v x)\,.
\end{align}
In the second diagram a gluon is first emitted and the incoming gluon is then absorbed. With the matrix
element~\eqref{m2},
\begin{align}\label{m2grr}
\ket{\big(\Psi^{a}_{i}(k)\big)^{2}_{g\rho\rho}}
=\ &\frac{1}{2}\int d^{3}p\,d^{3}q\,d^{3}r\,\ket{g^{b}_{j}(p)}\,
\frac{\bra{g^{b}_{j}(p)}H_{g}\ket{g^{c}_{l}(q)\,g^{d}_{m}(r)}\bra{g^{c}_{l}(q)\,g^{d}_{m}(r)}H_{g}\ket{g^{a}_{i}(k)}}
{\big(E_{g}(k)-E_{gg}(q,r)\big)\big(E_{g}(k)-E_{g}(p)\big)}\notag\\
=\ &-\frac{g^{2}}{4\pi^{3}}\int_{\Lambda/k^{+}}^{\vee/k^{+}}d\xi\int d^{2}\v p\,\ket{g^{b}_{j}(\xi k^{+},\v p)}\,
\frac{\sqrt{\xi}\;\v k^{i}\,\v p^{j}}{\v p^{2}\,(\xi\v k^{2}-\v p^{2})}\,\rho^{a}(\v k)\,\rho^{b}(-\v p)\,,
\end{align}
and, with Eq.~\eqref{i2},
\begin{align}\label{m2grrX}
\ket{\big(\Psi^{a}_{i}(k^{+},\v z')\big)^{2}_{g\rho\rho}}
=-\frac{g^{2}}{4\pi^{3}}\int_{\Lambda/k^{+}}^{\vee/k^{+}}d\xi\,\sqrt{\xi}\int_{\v x,\v y,\v z}\ket{g^{b}_{j}(\xi k^{+},\v z)}\,
\frac{X'^{i}\,Y^{j}}{(\xi Y^{2}-X'^{2})\,X'^{2}}\,\rho^{a}(\v x)\,\rho^{b}(\v y)\,.
\end{align}
The two diagrams order the charges differently. Their product is ambiguous up to the
commutator~\eqref{eq:rhoalgebra},
\begin{align}
[\rho^{a}(\v x),\rho^{b}(\v y)]=if^{abc}\rho^{c}(\v x)\,\delta^{(2)}(\v x-\v y)\,,\qquad
[\rho^{a}(\v k),\rho^{b}(-\v p)]=if^{abc}\rho^{c}(\v k-\v p)\,,
\end{align}
which reduces the number of valence degrees of freedom by one. We therefore write
\begin{align}
\rho^{a}(\v x)\,\rho^{b}(\v y)=\frac{1}{2}\{\rho^{a}(\v x),\rho^{b}(\v y)\}+\frac{1}{2}[\rho^{a}(\v x),\rho^{b}(\v y)]\,,
\end{align}
where the anticommutator depends on two valence degrees of freedom and is irreducible. The sum of the two
diagrams then splits as
\begin{align}
\ket{\big(\Psi^{a}_{i}(k)\big)^{1}_{g\rho\rho}}+\ket{\big(\Psi^{a}_{i}(k)\big)^{2}_{g\rho\rho}}
=\ket{\big(\Psi^{a}_{i}(k)\big)_{g\{\rho,\rho\}}}+\ket{\big(\Psi^{a}_{i}(k)\big)_{g[\rho,\rho]}}\,,
\end{align}
with
\begin{align}\label{onegrho3mom}
\ket{\big(\Psi^{a}_{i}(k)\big)_{g\{\rho,\rho\}}}
=\ &-\frac{g^{2}}{8\pi^{3}}\int_{\Lambda}^{\vee}\frac{dp^{+}}{\sqrt{k^{+}p^{+}}}\int d^{2}\v p\,\ket{g^{b}_{j}(p^{+},\v p)}\,
\frac{\v k^{i}\,\v p^{j}}{\v k^{2}\,\v p^{2}}\,\{\rho^{a}(\v k),\rho^{b}(-\v p)\}\notag\\
=\ &-\frac{g^{2}}{8\pi^{3}}\int_{\Lambda/k^{+}}^{\vee/k^{+}}\frac{d\xi}{\sqrt{\xi}}\int d^{2}\v p\,\ket{g^{b}_{j}(\xi k^{+},\v p)}\,
\frac{\v k^{i}\,\v p^{j}}{\v k^{2}\,\v p^{2}}\,\{\rho^{a}(\v k),\rho^{b}(-\v p)\}\,,\\
\ket{\big(\Psi^{a}_{i}(k)\big)_{g[\rho,\rho]}}
=\ &\frac{ig^{2}f^{abc}}{8\pi^{3}}\int_{\Lambda}^{\vee}\frac{dp^{+}}{\sqrt{k^{+}p^{+}}}\int d^{2}\v p\,\ket{g^{b}_{j}(p^{+},\v p)}\,
\frac{\v k^{i}\,\v p^{j}}{\v k^{2}\,\v p^{2}}\,\frac{k^{+}\v p^{2}+p^{+}\v k^{2}}{k^{+}\v p^{2}-p^{+}\v k^{2}}\,\rho^{c}(\v k-\v p)\notag\\
=\ &\frac{ig^{2}f^{abc}}{8\pi^{3}}\int_{\Lambda/k^{+}}^{\vee/k^{+}}\frac{d\xi}{\sqrt{\xi}}\int d^{2}\v p\,\ket{g^{b}_{j}(\xi k^{+},\v p)}\,
\frac{\v k^{i}\,\v p^{j}}{\v k^{2}\,\v p^{2}}\,\frac{\v p^{2}+\xi\v k^{2}}{\v p^{2}-\xi\v k^{2}}\,\rho^{c}(\v k-\v p)\,.\label{onegrho4mom}
\end{align}
The energy denominators cancel in the anticommutator part. It is the symmetrised product of the
leading-order absorption amplitude, Eq.~\eqref{eq:psirho}, and the leading-order emission amplitude,
Eq.~\eqref{lowav1mom}. In transverse coordinates,
\begin{align}\label{gtworhoX}
\ket{\big(\Psi^{a}_{i}(k^{+},\v z')\big)_{g\{\rho,\rho\}}}
=\ &-\frac{g^{2}}{8\pi^{3}}\int_{\Lambda/k^{+}}^{\vee/k^{+}}\frac{d\xi}{\sqrt{\xi}}\int_{\v x,\v y,\v z}
\ket{g^{b}_{j}(\xi k^{+},\v z)}\,\frac{X'^{i}\,Y^{j}}{X'^{2}\,Y^{2}}\,\{\rho^{a}(\v x),\rho^{b}(\v y)\}\,,\\
\ket{\big(\Psi^{a}_{i}(k^{+},\v z')\big)_{g[\rho,\rho]}}
=\ &\frac{ig^{2}f^{abc}}{8\pi^{3}}\int_{\Lambda/k^{+}}^{\vee/k^{+}}\frac{d\xi}{\sqrt{\xi}}\int_{\v x,\v z}
\ket{g^{b}_{j}(\xi k^{+},\v z)}\,\frac{X'^{i}\,X^{j}}{X'^{2}\,X^{2}}\,\frac{X'^{2}+\xi X^{2}}{X'^{2}-\xi X^{2}}\,\rho^{c}(\v x)\,.
\label{gtworhoX2}
\end{align}

\begin{figure}[t]
\centering
\begin{minipage}[b]{0.45\textwidth}\centering\tikzfig{gtovactog_tworho}\\[2pt](a)\end{minipage}\hfill
\begin{minipage}[b]{0.45\textwidth}\centering\tikzfig{gg_tworho}\\[2pt](b)\end{minipage}
\caption{One outgoing gluon, quadratic in $\rho$: (a) the incoming gluon is absorbed and a gluon is then
emitted; (b) a gluon is emitted and the incoming gluon is then absorbed.}
\label{fig:onegrhorho}
\end{figure}

%%%%%%%%%%%%%%%%%%%%%%%%%%%%%%%%%%%%%%%%%%%%%%%%%%%%%%%%%%%%%%%%%%%%%%%%%%%%%%%%
\paragraph{Two outgoing gluons, independent of $\rho$.}
The incoming gluon splits through the triple-gluon vertex, Fig.~\ref{fig:twog}(a). With the matrix
element~\eqref{m3},
\begin{align}\label{rhoindep}
\ket{\big(\Psi^{a}_{i}(k)\big)_{gg}}
=\ &\int d^{3}p\,d^{3}q\,\ket{g^{c}_{l}(q)\,g^{b}_{j}(p)}\,
\frac{\bra{g^{c}_{l}(q)\,g^{b}_{j}(p)}H_{ggg}\ket{g^{a}_{i}(k)}}{2\big(E_{g}(k)-E_{gg}(p,q)\big)}\notag\\
=\ &\int d^{3}p\,d^{3}q\,\ket{g^{c}_{l}(q)\,g^{b}_{j}(p)}\,
\frac{igf^{abc}\,\delta^{(3)}(k-p-q)}{8\pi^{3/2}\sqrt{k^{+}p^{+}q^{+}}\,
\Big(\frac{\v k^{2}}{k^{+}}-\frac{\v p^{2}}{p^{+}}-\frac{\v q^{2}}{q^{+}}\Big)}\notag\\
&\times\bigg[\Big(\v p^{i}-\v q^{i}+\frac{q^{+}-p^{+}}{k^{+}}\,\v k^{i}\Big)\delta_{jl}
+\Big(\v k^{j}+\v q^{j}-\frac{k^{+}+q^{+}}{p^{+}}\,\v p^{j}\Big)\delta_{il}
+\Big(\frac{k^{+}+p^{+}}{q^{+}}\,\v q^{l}-\v p^{l}-\v k^{l}\Big)\delta_{ij}\bigg]\notag\\
=\ &-\frac{igf^{abc}\sqrt{k^{+}}}{4\pi^{3/2}}\int_{\Lambda/k^{+}}^{1-\Lambda/k^{+}}d\xi\,\sqrt{\xi(1-\xi)}\int d^{2}\v p\notag\\
&\hspace{1cm}\times\ket{g^{c}_{l}\big((1-\xi)k^{+},\v k-\v p\big)\,g^{b}_{j}(\xi k^{+},\v p)}\,
\frac{\v p^{m}-\xi\v k^{m}}{(\v p-\xi\v k)^{2}}\,\mathcal{T}_{ijml}(\xi)\,.
\end{align}
The Fourier transform follows by shifting the integration variable $\v p$ and using Eq.~\eqref{i1}:
\begin{align}\label{gindeppos}
\ket{\big(\Psi^{a}_{i}(k^{+},\v z')\big)_{gg}}
=\ &-\frac{gf^{abc}\sqrt{k^{+}}}{4\pi^{3/2}}\int_{\Lambda/k^{+}}^{1-\Lambda/k^{+}}d\xi\,\sqrt{\frac{1-\xi}{\xi}}\int d^{2}z\,
\ket{g^{c}_{l}\big((1-\xi)k^{+},\v z\big)\,g^{b}_{j}\Big(\xi k^{+},\frac{\v z'-(1-\xi)\v z}{\xi}\Big)}\notag\\
&\times\frac{(z-z')^{m}}{(z-z')^{2}}\,\mathcal{T}_{ijml}(\xi)\,.
\end{align}
The second gluon sits at $(\v z'-(1-\xi)\v z)/\xi$, so that the plus-momentum-weighted centre of the pair
is at the position $\v z'$ of the incoming gluon.

\paragraph{Two outgoing gluons, linear in $\rho$.}
The valence charges emit a gluon while the incoming gluon is a spectator, Fig.~\ref{fig:twog}(b). With the
matrix element~\eqref{m2},
\begin{align}\label{eq:ggrhoone}
\ket{\big(\Psi^{a}_{i}(k)\big)_{gg\rho}}
=\ &\int d^{3}p\,d^{3}q\,\ket{g^{c}_{l}(q)\,g^{b}_{j}(p)}\,
\frac{\bra{g^{c}_{l}(q)\,g^{b}_{j}(p)}H_{g}\ket{g^{a}_{i}(k)}}{2\big(E_{g}(k)-E_{gg}(p,q)\big)}\notag\\
=\ &-\frac{g}{2\pi^{3/2}}\int_{\Lambda}^{\vee}\frac{dp^{+}}{\sqrt{p^{+}}}\int d^{2}\v p\,
\ket{g^{a}_{i}(k)\,g^{b}_{j}(p^{+},\v p)}\,\frac{\v p^{j}}{\v p^{2}}\,\rho^{b}(-\v p)\notag\\
=\ &-\frac{g\sqrt{k^{+}}}{2\pi^{3/2}}\int_{\Lambda/k^{+}}^{\vee/k^{+}}\frac{d\xi}{\sqrt{\xi}}\int d^{2}\v p\,
\ket{g^{a}_{i}(k)\,g^{b}_{j}(\xi k^{+},\v p)}\,\frac{\v p^{j}}{\v p^{2}}\,\rho^{b}(-\v p)\,.
\end{align}
In transverse coordinates,
\begin{align}\label{eq:ggrhooneX}
\ket{\big(\Psi^{a}_{i}(k^{+},\v z')\big)_{gg\rho}}
=\ &\frac{ig}{2\pi^{3/2}}\int_{\Lambda}^{\vee}\frac{dp^{+}}{\sqrt{p^{+}}}\int_{\v x,\v z}
\ket{g^{a}_{i}(k^{+},\v z')\,g^{b}_{j}(p^{+},\v z)}\,\frac{X^{j}}{X^{2}}\,\rho^{b}(\v x)\notag\\
=\ &\frac{ig\sqrt{k^{+}}}{2\pi^{3/2}}\int_{\Lambda/k^{+}}^{\vee/k^{+}}\frac{d\xi}{\sqrt{\xi}}\int_{\v x,\v z}
\ket{g^{a}_{i}(k^{+},\v z')\,g^{b}_{j}(\xi k^{+},\v z)}\,\frac{X^{j}}{X^{2}}\,\rho^{b}(\v x)\,.
\end{align}
The amplitude is that of the leading-order emission from the vacuum, Eq.~\eqref{lowav1mom}, as required by
the $\mf{A}_{2}$ term of Eq.~\eqref{eq:Omegag}.

\begin{figure}[t]
\centering
\begin{minipage}[b]{0.45\textwidth}\centering\tikzfig{g_togg}\\[2pt](a)\end{minipage}\hfill
\begin{minipage}[b]{0.45\textwidth}\centering\tikzfig{gtog_plusg}\\[2pt](b)\end{minipage}
\caption{Two outgoing gluons: (a) splitting of the incoming gluon, independent of $\rho$; (b) emission
by the valence charges with the incoming gluon as a spectator.}
\label{fig:twog}
\end{figure}

%%%%%%%%%%%%%%%%%%%%%%%%%%%%%%%%%%%%%%%%%%%%%%%%%%%%%%%%%%%%%%%%%%%%%%%%%%%%%%%%
\paragraph{Normalisation.}
To order $g^{2}$ the dressed one-gluon state is the sum of the contributions above:
\begin{align}\label{Eq:intergluon}
\ket{\Psi^{a}_{i}(k)}=\ &\mf{N}\ket{g^{a}_{i}(k)}\notag\\
&+\ket{\big(\Psi^{a}_{i}(k)\big)_{\rho}}+\ket{\big(\Psi^{a}_{i}(k)\big)_{gg}}+\ket{\big(\Psi^{a}_{i}(k)\big)_{gg\rho}}\notag\\
&+\ket{\big(\Psi^{a}_{i}(k)\big)^{1}_{g\rho,\,{\rm inst}}}+\ket{\big(\Psi^{a}_{i}(k)\big)^{2}_{g\rho,\,{\rm inst}}}
+\ket{\big(\Psi^{a}_{i}(k)\big)^{1}_{g\rho}}+\ket{\big(\Psi^{a}_{i}(k)\big)^{2}_{g\rho}}\notag\\
&+\ket{\big(\Psi^{a}_{i}(k)\big)^{1}_{g\rho\rho}}+\ket{\big(\Psi^{a}_{i}(k)\big)^{2}_{g\rho\rho}}\,.
\end{align}
The first line is proportional to the free gluon state, the second line contains the contributions of
order $g$, and the last two lines those of order $g^{2}$. The coefficient $\mf{N}$ is fixed by the
normalisation of the dressed state,
\begin{align}
\braket{\Psi^{a'}_{i'}(k')|\Psi^{a}_{i}(k)}=\delta^{aa'}\,\delta_{ii'}\,\delta^{(3)}(k-k')\,.
\end{align}
We need $\mf{N}$ only to order $g^{2}$, so few overlaps contribute: products of two order-$g^{2}$ terms,
or of an order-$g$ and an order-$g^{2}$ term, are of higher order. The following combinations vanish:
\begin{align}
&\braket{g^{a'}_{i'}(k')|\big(\Psi^{a}_{i}(k)\big)^{1}_{g\rho,\,{\rm inst}}}
+\braket{g^{a'}_{i'}(k')|\big(\Psi^{a}_{i}(k)\big)^{2}_{g\rho,\,{\rm inst}}}+{\rm h.c.}=0\,,\notag\\
&\braket{g^{a'}_{i'}(k')|\big(\Psi^{a}_{i}(k)\big)^{1}_{g\rho}}
+\braket{g^{a'}_{i'}(k')|\big(\Psi^{a}_{i}(k)\big)^{2}_{g\rho}}+{\rm h.c.}=0\,,\notag\\
&\braket{g^{a'}_{i'}(k')|\big(\Psi^{a}_{i}(k)\big)^{1}_{g\rho\rho}}+{\rm h.c.}
+\braket{\big(\Psi^{a'}_{i'}(k')\big)_{\rho}|\big(\Psi^{a}_{i}(k)\big)_{\rho}}=0\,,\notag\\
&\braket{\big(\Psi^{a'}_{i'}(k')\big)_{gg\rho}|\big(\Psi^{a}_{i}(k)\big)_{gg}}+{\rm h.c.}=0\,,
\end{align}
where h.c.\ denotes the overlap with bra and ket exchanged and $k\leftrightarrow k'$. The third line
cancels exactly, including the order of the charges: both terms are proportional to
$\rho^{a'}(-\v k')\rho^{a}(\v k)$. The splitting term gives
\begin{align}
\braket{\big(\Psi^{a'}_{i'}(k')\big)_{gg}|\big(\Psi^{a}_{i}(k)\big)_{gg}}
=-\delta^{aa'}\,\delta_{ii'}\,\delta^{(3)}(k-k')\,\frac{g^{2}N_{c}}{8\pi^{3}}
\Big(2\ln\frac{\Lambda}{k^{+}}+\frac{11}{6}\Big)\int\frac{d^{2}\tilde{\v p}}{\tilde{\v p}^{2}}=0\,,
\end{align}
since the transverse integral is scaleless and vanishes in dimensional regularisation. The remaining
terms are
\begin{align}
&\braket{g^{a'}_{i'}(k')|\big(\Psi^{a}_{i}(k)\big)^{2}_{g\rho\rho}}+{\rm h.c.}
=-\frac{g^{2}}{4\pi^{3}\sqrt{k^{+}k'^{+}}}\,\frac{\v k^{i}\,\v k'^{i'}}{\v k^{2}\,\v k'^{2}}\,
\rho^{a}(\v k)\,\rho^{a'}(-\v k')\,,\notag\\
&\braket{\big(\Psi^{a'}_{i'}(k')\big)_{gg\rho}|\big(\Psi^{a}_{i}(k)\big)_{gg\rho}}
=\frac{g^{2}}{4\pi^{3}\sqrt{k^{+}k'^{+}}}\,\frac{\v k^{i}\,\v k'^{i'}}{\v k^{2}\,\v k'^{2}}\,
\rho^{a}(\v k)\,\rho^{a'}(-\v k')\notag\\
&\hspace{4.6cm}+\delta^{aa'}\,\delta_{ii'}\,\delta^{(3)}(k-k')\,\frac{g^{2}}{4\pi^{3}}
\int_{\Lambda}^{\vee}\frac{dp^{+}}{p^{+}}\int d^{2}\v p\,\frac{\rho^{b}(-\v p)\,\rho^{b}(\v p)}{\v p^{2}}\,.
\end{align}
The first terms cancel, again with the same order of the charges. The second term of the last line is
the exchange term of the two identical gluons. The normalisation condition then gives
\begin{align}\label{eq:Nonegluon}
\mf{N}=1-\frac{g^{2}}{8\pi^{3}}\ln\Big(\frac{\vee}{\Lambda}\Big)\int d^{2}\v p\,\frac{\rho^{a}(-\v p)\,\rho^{a}(\v p)}{\v p^{2}}\,.
\end{align}
This is the same operator as for the dressed vacuum, Eqs.~\eqref{eq:Nmom} and~\eqref{eq:Npos}. It must
be, since $\Omega$ is a single operator and $\mf{N}$ multiplies both $\ket{0}$ in Eq.~\eqref{eq:Omegavac}
and $\ket{g^{a}_{i}(k)}$ in Eq.~\eqref{eq:Omegag}.

Comparing Eq.~\eqref{Eq:intergluon} with Eq.~\eqref{eq:Omegag} identifies the coefficients:
$\ket{(\Psi^{a}_{i}(k))_{\rho}}$ gives $\mf{A}^{(1)}_{1}$, the one-gluon components of order $g^{2}$ give
$\mf{B}_{3}$, $\ket{(\Psi^{a}_{i}(k))_{gg}}$ gives $\mf{C}_{1}$, and $\ket{(\Psi^{a}_{i}(k))_{gg\rho}}$
reproduces the $\mf{A}^{(1)}_{2}$ term.
